% One complete cyber-response case, 135--165 minutes.
\section{Cyber response}

\begin{frame}[t]{Cyber response: problem, workflow and comparison}
\noindent
{\centering
\begin{tikzpicture}[every node/.style={font=\small,align=center},
 gridstage/.style={draw=forest,rounded corners,fill=black!3,minimum height=.9cm,text width=2.65cm}]
 \node[gridstage] (a) at (0,0) {Physical grid\\\textbf{unknown state}};
 \node[gridstage] (b) at (3.25,0) {Sensors\\\textbf{measurements}};
 \node[gridstage] (c) at (6.5,0) {Estimator\\\textbf{estimated state}};
 \node[gridstage] (d) at (9.75,0) {Controller\\\textbf{proposed action}};
 \draw[->,thick,forest] (a)--(b); \draw[->,thick,forest] (b)--(c); \draw[->,thick,forest] (c)--(d);
 \node[above=.35cm,font=\small,text=orange] (attack) at (b.north) {Possible corruption};
 \draw[->,thick,orange] (attack)--(b);
\end{tikzpicture}\par}
\operatorneed{Which evidence can Planning rely on, and what must Verification reject before an operator reviews the action?}
\small\textbf{Case sequence:} residual tests (C1--C2), trusted-data recovery (C3), initial supervision (C4), complete refinement for two attacks (P6). The numerical model has four DC buses, three unknown angles and eight original measurements.
\sources{State estimation: Schweppe and Wildes, 1970, \url{https://doi.org/10.1109/TPAS.1970.292678}. Trust is an explicit tutorial assumption.}
\note{135--137 minutes. Connect directly to the previous question: a model can validate a decision against the wrong estimated state. The four-bus example is separate from the Greensboro feeder. Measurements comprise four branch flows, three injections, and one angle. The three unknowns are non-slack voltage angles in radians. Synthetic truth is known to the experimenter, enabling evaluation, but an actual operator generally does not know it. No full-feeder attack detector is claimed.}
\end{frame}

\begin{frame}[t]{Fit the readings, then examine what the fit leaves unexplained}
\small
\textbf{Weighted least squares (WLS)} fits a physical measurement model:
\[
z=Hx+\epsilon,\qquad \epsilon\sim\mathcal N(0,R),\qquad
\hat x=\arg\min_{x\in\mathbb R^3}(z-Hx)^\top R^{-1}(z-Hx).
\]
\footnotesize
\renewcommand{\arraystretch}{1.02}
\begin{tabular}{@{}>{\raggedright\arraybackslash}p{.22\linewidth}>{\raggedright\arraybackslash}p{.70\linewidth}@{}}
\toprule Quantity & Meaning and units in this four-bus DC example \\
\midrule
$x$, $\hat x$ & True and estimated three non-reference voltage angles (rad). Reference angle is fixed at bus 0. \\
$z\in\mathbb R^8$ & Four branch flows and three injections (p.u.), plus one angle (rad); mixed units are handled by weighting. \\
$H\in\mathbb R^{8\times3}$ & Linear measurement sensitivity; each row maps radians into that sensor's units. Full column rank 3 is required. \\
$\epsilon$, $R$ & Measurement noise and its covariance; here independent Gaussian noise with $R=\operatorname{diag}(\sigma_i^2)$; $\sigma_i$ is sensor standard deviation in its own units. \\
WLS objective & Dimensionless normalized squared mismatch; smaller noise variance gives a reading greater weight. \\
\bottomrule
\end{tabular}
\smallskip\par\textbf{Interpretation:} good agreement indicates consistency with the model/noise assumptions, not trustworthy measurements. The code checks observability and solves the weighted system without explicitly forming the normal-matrix inverse.
\sources{Schweppe and Wildes, 1970: \url{https://doi.org/10.1109/TPAS.1970.292678}.}
\note{137--140 minutes. Let the audience answer before the reveal. Explain z as the eight measurements, H as their linear dependence on the three angles, and R as the declared noise covariance. The implementation checks full column rank and uses weighted least squares, rather than explicitly inverting the normal matrix. An isolated error may fail consistency, while a coordinated error may not. The next checkpoint quantifies this distinction; the following plot makes the failure visible.}
\end{frame}

\begin{frame}[t]{From residuals to an alarm and a suspect measurement}
\small
After WLS, define residual $r=z-H\hat x$ in the measurements' units. The dimensionless global statistic is
\[
J=r^\top R^{-1}r,\qquad J>\chi^2_{\nu,0.99}=15.0863,\qquad \nu=m-n=8-3=5.
\]
\footnotesize $z$ is the measurement vector, $H$ the linear sensitivity, $\hat x$ the angle estimate, and $R$ the stated noise covariance. $m$ counts measurements; $n$ counts observable states. $\chi^2_{\nu,0.99}$ is the 99th percentile: a nominal 1\% false-alarm probability under the correct linear model and Gaussian noise, not a measured rate here.
\[
\Omega=R-H(H^\top R^{-1}H)^{-1}H^\top,\qquad
r_i^N=\frac{|r_i|}{\sqrt{\Omega_{ii}}}\quad(\Omega_{ii}>0).
\]
$\Omega$ is residual covariance; $i$ is a measurement index, $r_i$ its residual, and $r_i^N$ its dimensionless normalized magnitude. The largest value localizes a suspect reading; it is not a guaranteed attacker identifier. A zero residual variance makes this score undefined.
\smallskip\par\small
\textbf{Before recovery:} investigate inconsistency and recheck observability after removal. Topology/model error, sensor failure and malicious corruption can all produce an alarm.
\sources{State estimation: \href{https://doi.org/10.1109/TPAS.1970.292678}{Schweppe and Wildes (1970)}; bad-data residuals: \href{https://doi.org/10.1109/T-PAS.1975.31858}{Handschin et al. (1975)}. Exact calculation: \texttt{weighted\_least\_squares}.}
\note{About one minute within the existing 140--143-minute detector explanation. The covariance expression explains why dividing residuals by raw sensor standard deviations is not the same localization statistic. Assumptions matter: the nominal chi-square calibration is exact for this linear full-rank Gaussian construction, not a universal guarantee for a nonlinear AC estimator. The code flags undefined residual normalizations and checks rank after removal.}
\end{frame}

\begin{frame}[t]{C1: use one detector on three measurement cases}
\checkpoint{C1}{Run nominal noise, one biased reading, and a coordinated stealth attack.}
\cellall{C1}
\observethis{Compare the residual alarm with the actual control outcome. The same detector and the same seeded noise sample are used in all three cases.}
\footnotesize $J=r^\top R^{-1}r$ is dimensionless, with residual $r=z-H\hat x$; $z$ denotes measurements, $H$ their linear model, $\hat x$ the angle estimate and $R$ the noise covariance. Alarm threshold: $15.0863$ for five residual degrees of freedom.\par
\smallskip Isolated corruption adds 0.08 p.u. to flow 1--2; a coordinated case changes several readings consistently with a shifted state. Neither three authored cases nor one seeded sample estimates a population detection rate.
\sources{Nominal 1\% test assumes the stated Gaussian noise and correct linear model; Handschin et al., 1975: \url{https://doi.org/10.1109/T-PAS.1975.31858}.}
\note{140--143 minutes. The chi-square threshold uses five residual degrees of freedom and alpha 0.01. This is a nominal calibration, not a measured false-alarm rate from three trials. All cases use the same model and fixed noise sample. The isolated corruption adds 0.08 p.u. to flow 1--2, measurement index 1. Its largest normalized residual identifies that measurement in this fixture, and control is held. The stealth case uses a=Hc with c=(0.01,0.025,0) radians. Detailed residual covariance and localization equations are in the appendix.}
\end{frame}

\begin{frame}[t]{C2: a small residual can accompany a large state error}
\checkpoint{C2}{Run the plot, then compare the two vertical axes.}
\cellcode{C2}{1}{1}
{\centering\includegraphics[width=.92\linewidth,height=.40\textheight,keepaspectratio]{\figfile{workshop_cyber.png}}\par}
\textcolor{forest}{\textbf{Question:}} Why do the nominal and stealth residual statistics both equal 3.9202?\par
\uncover<2->{\textbf{Reason:} the corruption $a=Hc$ shifts the estimated angles by $c$ while leaving the residual unchanged.}
\par\footnotesize $a$ is the eight-value added measurement vector (sensor units); $H$ is the linear sensitivity; $c=(0.01,0.025,0)$ rad is the three-angle shift. Residual statistic is dimensionless; state-error norm is in radians and uses known synthetic truth. The attack requires access to the affected channels and knowledge of $H$.
\sources{Stealth construction: Liu et al., 2009, \url{https://doi.org/10.1145/1653662.1653666}. Numerical evidence: \texttt{results/standard\_formulations.json}.}
\note{143--146 minutes. Pause before the reason. The isolated statistic is 277.4974, above the 15.0863 threshold. Nominal and stealth statistics are both 3.9202, below it. Their state error norms differ: 0.000237 versus 0.026780 radians. The full identity is z'=z+Hc, xhat'=xhat+c, hence z'-H xhat'=r. The stealth construction assumes knowledge of H and access to the unprotected channels required by a. It is a designed negative control, not a detection-rate survey. The right-hand state error uses known synthetic truth.}
\end{frame}

\begin{frame}[t]{Recovery needs new evidence before it needs a new action}
\footnotesize
\textbf{Evidence assumption:} three additional protected angles give $H_{\mathcal T}=I_3$ (the $3\times3$ identity). This set is observable with \textbf{zero residual redundancy}; it is not a subset of the original eight readings. Sensor trust is assumed, not authenticated by WLS or Llama.
\setlength{\abovedisplayskip}{4pt}\setlength{\belowdisplayskip}{4pt}\[
\min_{0\leq u\leq0.5}\ \tfrac12u^2
\quad\text{s.t.}\quad
|f_{\rm pred}(u)|=|F\hat x_{\mathcal T}+b u|\leq\overline f-0.01.
\]
\footnotesize
\renewcommand{\arraystretch}{1.0}
\begin{tabular}{@{}>{\raggedright\arraybackslash}p{.22\linewidth}>{\raggedright\arraybackslash}p{.70\linewidth}@{}}
\toprule Symbol & Physical meaning / units \\
\midrule
$u$ & Added active injection at node 2, balanced at the source; scalar in $[0,0.5]$ p.u. Minimize its squared effort. \\
$\hat x_{\mathcal T}$ & Three angles estimated from protected readings (rad); their stated noise remains present. \\
$F$, $b$ & DC angle-to-flow matrix (p.u./rad) and four-branch injection sensitivity (p.u. flow per p.u. injection). \\
$f_{\rm pred}$, $\overline f$ & Predicted signed flows and absolute limits (p.u.). Limits: $(0.35,0.18,0.35,0.35)$, applied branch by branch. \\
$0.01$ & Planning margin (p.u.), a declared engineering choice rather than a calibrated uncertainty bound. \\
\bottomrule
\end{tabular}
\par\textbf{Check consequences:} validate both injections against the same synthetic truth. This example includes only DC active-power constraints.
\sources{Stealth limitation: Liu, Ning, and Reiter, 2009: \url{https://doi.org/10.1145/1653662.1653666}.}
\note{146--150 minutes. Say: The residual test cannot distinguish these two states. What additional information could break that ambiguity? Introduce the explicitly protected angle measurements with H sub T equal to I3. Full rank identifies the state, while three measurements for three unknowns gives zero residual degrees of freedom; there is no chi-square detection redundancy in this protected set. The numerical experiment assumes trust and does not implement authentication. The scalar control uses the DC flow map F and injection sensitivity F B inverse d2. Bounds, finite values, and predicted constraints are checked after the solve. The 0.01 p.u. margin is an engineering choice, not a calibrated uncertainty guarantee. This is a small DC control problem, not AC feeder redispatch.}
\end{frame}

\begin{frame}[t]{C3: compare the actions, then evaluate their consequences}
\checkpoint{C3}{Lines 1--4 compare control from the stealth and protected estimates.}
\cellcode{C3}{1}{4}
{\centering
\begin{tabular}{lrr}
\toprule
Estimate used & Injection $u$ (p.u.) & Actual flow $1$--$2$ (p.u.) \\
\midrule
Stealth-corrupted & 0 & 0.24000 \\
Protected measurements & 0.12169 & 0.17046 \\
\bottomrule
\end{tabular}\par}
\observethis{Limit: \textbf{0.18 p.u.} Only protected-data control removes this toy overload.}
\footnotesize $u$ is added active injection at node 2, balanced at the source. Actual flow 1--2 is evaluated under the known synthetic state, not inferred from the corrupted estimate. All powers share the toy model's arbitrary fixed per-unit base.
\sources{Estimation context: Schweppe and Wildes, 1970, \url{https://doi.org/10.1109/TPAS.1970.292678}. Control/evaluation: tutorial four-bus DC experiment.}
\note{150--153 minutes. The corrupted estimate predicts a feasible state and chooses u=0, but actual toy flow remains 0.24 p.u. The protected estimate chooses u=0.1216936673 p.u. Its predicted flow is 0.17 p.u.; evaluation against synthetic truth gives 0.1704607616 p.u., below 0.18. The small prediction/evaluation difference is due to protected measurement noise. Explain that the experiment evaluates both actions under the same known physical toy state, which prevents a favorable result merely because an estimate reports its own constraints as feasible.}
\end{frame}

\begin{frame}[t]{C3: check every branch against its own limit}
\checkpoint{C3}{Lines 5--6 plot the consequence and report estimator redundancy.}
\cellcode{C3}{5}{6}
{\centering\includegraphics[width=.94\linewidth,height=.44\textheight,keepaspectratio]{\figfile{workshop_trusted_control.png}}\par}
\textbf{Check the consequence:} compare every bar with its branch limit.
\sources{Stealth limitation: Liu et al., 2009, \url{https://doi.org/10.1145/1653662.1653666}. Post-control flows: \texttt{results/standard\_formulations.json}.}
\note{153--156 minutes. Point to flow 1--2 above its limit under the stealth action and below it after protected-data redispatch. Then check the other branches. The branch limits are 0.35, 0.18, 0.35, 0.35 p.u. The chart uses known synthetic truth, not field measurements. The result omits voltage magnitudes, unbalance, reactive capability, protection, and communication authentication. The printed protected-set residual degrees of freedom are zero; this is not a detector success claim for that set.}
\end{frame}

\begin{frame}[t]{C4: the agent can explain evidence, but cannot establish trust}
\deepdive
\checkpoint{C4}{Inspect the recorded actual Llama response for estimation and recovery.}
\cellall{C4}
{\centering
\begin{tikzpicture}[every node/.style={font=\small,align=center},
 cresponse/.style={draw=forest,rounded corners,text width=2.7cm,minimum height=1.05cm}]
 \node[cresponse] (a) at (0,0) {\textbf{Planning}\\Llama selects\\WLS/control tool};
 \node[cresponse] (b) at (3.35,0) {\textbf{Verification}\\DC checks PASS\\Trust is assumed};
 \node[cresponse] (c) at (6.7,0) {\textbf{Explanation}\\Llama: REVIEW};
 \node[cresponse] (d) at (10.05,0) {\textbf{Human review}\\Pending: request\\further evidence};
 \draw[->,thick,forest] (a)--(b); \draw[->,thick,forest] (b)--(c); \draw[->,thick,forest] (c)--(d);
\end{tikzpicture}\par}
\observethis{The fixed route recommends the same numerical result. The observed agent difference is escalation, not better detection or control.}
\sources{Tool/observation pattern: Yao et al., 2023, \url{https://arxiv.org/abs/2210.03629}. Actual responses: \texttt{results/standard\_agent\_workflow.json}.}
\note{156--159 minutes. This is a recorded trace, not a new Llama call. Planning call 7 selects WLS\_TRUSTED\_REDISPATCH. Explanation call 8 selects REVIEW despite the passed numerical checks. The fixed comparator returns MODEL\_RECOMMEND. The control and detector are identical, so this trial establishes no agent advantage in detection or control quality. The wrapper explicitly supplies the trusted-data branch; it does not discover which sensors deserve trust. Separate the deterministic numerical metric card from the actual model's qualitative response. This visible separation is a useful design property even when the agent escalates unnecessarily.}
\end{frame}

\begin{frame}[t]{Your decision: what evidence justifies recovery?}
\exercise{3}{Explain why the stealth estimate passes the residual test but gives a poor action. What does the protected set add? What fails if that set loses rank? Why is \textit{no alarm} insufficient?}
\answer{A model-consistent attack can move the estimated state without moving the residual. Trusted, observable data can identify the state; rank loss leaves it underdetermined. Without justified evidence, escalate for review.}
\sources{Stealth limitation: Liu et al., 2009, \url{https://doi.org/10.1145/1653662.1653666}. Observability: Schweppe and Wildes, 1970, \url{https://doi.org/10.1109/TPAS.1970.292678}.}
\note{159--162 minutes. Reserve three full minutes before revealing the answer. This follows C4, matching the notebook exercise. Expected explanation: a=Hc preserves the residual while changing the estimated angles; the chosen u=0 then leaves an overload under the true toy state. The protected H sub T=I3 is full rank, but has zero residual redundancy. Rank deficiency is rejected by the implementation. No residual alarm is a consistency statement, not proof of no attack. If the trusted evidence is missing or unobservable, the framework must not invent a state or control. Holding a previous setpoint is not automatically safe after an unknown physical change.}
\end{frame}

% Historical standard-method trace; one case-specific evidence frame.
\begin{frame}[t]{C4: the actual cyber agent retains a provisional candidate}
\deepdive
\checkpoint{C4}{Recorded four-bus DC workflow; trusted measurements are supplied as an assumption.}
\footnotesize
\begin{tabular}{@{}p{.23\linewidth}p{.43\linewidth}p{.28\linewidth}@{}}
\toprule Stage / method & Actual response & Next action \\
\midrule
Planning / Llama & \texttt{WLS\_TRUSTED\_REDISPATCH} & Run the permitted tool \\
Estimation / trusted WLS & Rank 3; three angles; residual degrees of freedom 0 & Form the DC control problem \\
Proposal / redispatch & $u=0.12169$ p.u. flexible injection & Evaluate under synthetic truth \\
Verification / DC checks & Flow $1$--$2=0.17046<0.18$ p.u.; zero violations & Prepare numerical report \\
Explanation / Llama & \texttt{REVIEW}, despite passed modeled checks & Escalate; candidate provisional \\
Human review / authority & Acknowledgment remains pending & No equipment command \\
\bottomrule
\end{tabular}
\par\smallskip
\textbf{Actual Planner reason:} ``This tool is authorized for the given use case and scope, and it is designed for a specific DC redispatch example with trusted measurements.''
\par\smallskip
\textbf{Strong baseline:} fixed supervision uses the same $u$, has the same flows and returns \texttt{MODEL\_RECOMMEND}. No added Llama detection or control benefit is demonstrated; its escalation remains visible.
\sources{Exact responses: \texttt{results/standard\_agent\_workflow.json}, cyber calls 7--8. Tool evidence: \texttt{results/standard\_formulations.json}. Stealth limitation: Liu et al., \url{https://doi.org/10.1145/1653662.1653666}.}
\note{Read the returned enum, not what we hoped the agent would do. Llama chooses the correct tool family on Planning call7. The numerical solver produces injection0.12169366726990237 per unit, and simulation using known synthetic truth returns branch1--2 flow0.1704607615600558 per unit. All four modeled branch limits pass. On Explanation call8, the model chooses REVIEW even though its reason says the workflow completed successfully and passed modeled checks. The implementation therefore returns REVIEW; the numerical candidate must remain provisional. Do not silently convert the narrative into MODEL\_RECOMMEND. The fixed adaptive comparator uses the same protected readings and numerical method, and gets the same action. Compared with control from the corrupted stealth estimate, protected-data recovery lowers the actual toy flow from0.24 to0.17046 and removes one modeled overload. That improvement belongs to the additional evidence and redispatch, not to Llama. Full-rank trusted measurements make the toy state identifiable, but they have zero residual redundancy and their authenticity is assumed, not established by the framework. This is a recorded standard-method trace from September21, separate from the current continuous voltage study.}
\end{frame}

% Twelve optional code-detail frames. Exact source excerpts; no new experiments.
% Insert after C4 evidence and before the final cyber comparison.
\begingroup
% Compact exact-source listings locally; subsequent deck sections retain their style.
\renewcommand{\implcode}[3]{\lstinputlisting[style=workshoppython,basicstyle=\ttfamily\fontsize{7}{8}\selectfont,firstline=#2,lastline=#3]{snippets/impl_#1.py}}

\begin{frame}[t]{Cyber code 1/12: construct a measurement problem we can audit}
\deepdive
\checkpoint{C1}{The notebook calls this four-bus DC fixture; it is separate from Greensboro.}
\implcode{standard_formulations}{420}{428}
\footnotesize
\textbf{Read the arrays:} \texttt{incidence} links four branches to three non-reference angles. \texttt{F} maps angles to branch flows; \texttt{B} maps angles to bus injections. Stacking four flow rows, three injection rows and one angle row gives \texttt{H.shape=(8,3)}.
\par\smallskip
\textbf{Units and observed output:} angles are radians; powers use an arbitrary fixed per-unit base. Sensor standard deviations are 0.004 p.u. for flows, 0.006 p.u. for injections and 0.0005 rad for the angle. Seed 11 makes the shared noise sample reproducible.
\sources{Exact implementation: \texttt{standard\_formulations.py:420--428}; recorded evidence: \texttt{results/standard\_formulations.json}, \texttt{cyber.measurement\_model}. WLS: \href{https://doi.org/10.1109/TPAS.1970.292678}{Schweppe and Wildes (1970)}.}
\note{Optional code detail. Read C1 in the notebook first, then locate the function it invokes. The Python variable truth is known to this synthetic experiment, not to a real control-room operator. H is a fixed DC measurement sensitivity; no learned model or Llama calculation creates it. Explain that the three voltage angles omit the reference angle, making an identifiable estimate possible. The factor multiplying each incidence row is the declared branch coefficient. All three attack scenarios reuse this network and noise realization, so a changed output is attributable to changed measurements within this fixture.}
\end{frame}

\begin{frame}[t]{Cyber code 2/12: establish observability before fitting a state}
\deepdive
\checkpoint{C1}{Inspect the estimator called for every measurement case.}
\implcode{standard_formulations}{373}{380}
\footnotesize
\textbf{Why these operations matter:} full column rank means the measurements distinguish the three modeled angles. Dividing each row and reading by its standard deviation normalizes mixed sensor units before least squares.
\par\smallskip
\textbf{Actual nominal case:} rank = 3; angle-error norm = 0.00023743 rad. \texttt{estimate} is the fitted angle vector; \texttt{residual} is the remaining measurement mismatch. \texttt{covariance} and \texttt{omega} quantify estimate and residual covariance under the stated linear Gaussian assumptions.
\sources{Source: \texttt{standard\_formulations.py:373--380}; output: \texttt{cyber.cases[0]} in \texttt{results/standard\_formulations.json}. \href{https://doi.org/10.1109/TPAS.1970.292678}{Schweppe and Wildes (1970)}.}
\note{Optional code detail. Rank deficiency raises an exception rather than producing an unidentifiable control state. The implementation uses numpy.linalg.lstsq for the estimate, rather than explicitly forming the normal-equation inverse. The covariance uses a linear solve. The nominal angle error is computed against synthetic truth for evaluation; it is not available as an online diagnostic in a real system. A full-rank estimator can still fit malicious but physically consistent data, which motivates the next checks.}
\end{frame}

\begin{frame}[t]{Cyber code 3/12: derive the residual test from its assumptions}
\deepdive
\checkpoint{C1}{Trace the statistic and threshold printed by the notebook.}
\implcode{standard_formulations}{380}{387}
\footnotesize
\textbf{Global decision:} \texttt{statistic} is the sum of squared residuals normalized by sensor variance. With eight measurements and rank three, \texttt{dof=5}; \texttt{alpha=.01} gives threshold 15.08627247. The nominal statistic 3.92022169 produces no alarm.
\par\smallskip
\textbf{Localization is different:} \texttt{normalized} divides each residual by its own residual standard deviation. Zero residual variance is left undefined. The largest value identifies a suspect measurement in this example, not a guaranteed attacker.
\sources{Source: \texttt{standard\_formulations.py:380--387}; outputs: \texttt{cyber.cases[*].chi\_square}. Residual tests: \href{https://doi.org/10.1109/T-PAS.1975.31858}{Handschin et al. (1975)}.}
\note{Optional code detail. Explain why rank, rather than simply the number of state columns, determines the residual degrees of freedom; this implementation already requires full column rank. Alpha is a declared nominal false-alarm probability under the correct linear model and known independent Gaussian noise. Three designed cases do not estimate an empirical false-alarm or detection rate. The finite diagonal test prevents division by a residual variance too close to zero.}
\end{frame}

\begin{frame}[t]{Cyber code 4/12: an isolated bad reading stops the control route}
\deepdive
\checkpoint{C1}{Follow the deliberately corrupted reading through the common estimator.}
\implcode{standard_formulations}{429}{436}
\footnotesize
\textbf{Actual isolated result:} adding 0.08 p.u. to measurement 1 (flow 1--2) raises the residual statistic to 277.49745, above 15.08627. The branch returns \texttt{HOLD}; no redispatch is proposed from that estimate.
\par\smallskip
\textbf{Useful automation:} the deterministic test prevents a downstream controller from silently using a rejected estimate. Neither this alarm nor its successful localization identifies malicious intent; measurement and model errors also cause inconsistency.
\sources{Source: \texttt{standard\_formulations.py:429--436}; output: \texttt{cyber.cases[1]}, \texttt{results/standard\_formulations.json}. \href{https://doi.org/10.1109/T-PAS.1975.31858}{Handschin et al. (1975)}.}
\note{Optional code detail. The loop applies one estimator and one alarm rule to the nominal, isolated and stealth measurements. The highlighted conditional is conventional deterministic control logic, not an LLM judgment. HOLD means this routine does not propose a control from the suspect estimate. It does not prove that the current physical operating point remains safe. Index 1 is zero-based and identifies the second measurement, flow 1--2.}
\end{frame}

\begin{frame}[t]{Cyber code 5/12: coordinated corruption can leave the test unchanged}
\deepdive
\checkpoint{C2}{Connect the unchanged residual to the attack construction, not to a plot artifact.}
\implcode{standard_formulations}{426}{433}
\footnotesize
\textbf{Mechanism:} \texttt{H @ shift} adds measurements consistent with an angle shift of $(0.01,0.025,0)$ rad. WLS absorbs the shift into the estimate while leaving its residual unchanged.
\par\smallskip
\textbf{Actual negative control:} nominal and stealth residual statistics both equal 3.92022. The angle-error norm rises from 0.00023743 to 0.026779996 rad. The stealth estimate chooses zero injection, leaving actual flow 1--2 at 0.24 p.u., above its 0.18 limit.
\sources{Source: \texttt{standard\_formulations.py:426--433}; output: \texttt{cyber.cases[2]}, \texttt{results/standard\_formulations.json}. False-data injection: \href{https://doi.org/10.1145/1653662.1653666}{Liu et al. (2009)}.}
\note{Optional code detail. This is a designed Hc attack, where H is the measurement sensitivity and c is the selected angle shift. It assumes knowledge of H and access to the required unprotected channels; it is not an arbitrary sparse-attack detection result. A small residual checks consistency, not authenticity. The true-state error and actual branch violation are measurable here only because the fixture supplies known synthetic truth. The failure exists before any language-model call.}
\end{frame}

\begin{frame}[t]{Cyber code 6/12: removal and protected evidence solve different problems}
\deepdive
\checkpoint{C3}{The recovery branch receives additional protected angle readings explicitly.}
\implcode{standard_formulations}{439}{446}
\footnotesize
\textbf{Isolated-error branch:} remove suspect index 1, then refit. Rank remains 3; seven measurements leave four residual degrees of freedom. The statistic falls to 2.32189, below its new threshold 13.27670.
\par\smallskip
\textbf{Protected-data branch:} three additional angles use the identity sensitivity and standard deviation 0.00005 rad. They are not a subset discovered by the LLM. Rank = 3, residual degrees of freedom = 0; the chi-square alarm is undefined, not a successful attack-detection result.
\sources{Source: \texttt{standard\_formulations.py:439--446}; outputs: \texttt{cyber.isolated\_removal}, \texttt{cyber.trusted\_recovery}. Observability: \href{https://doi.org/10.1109/TPAS.1970.292678}{Schweppe and Wildes (1970)}.}
\note{Optional code detail. Do not confuse deletion of an isolated suspect with recovery from a model-consistent coordinated attack. The former retains redundant measurements and recomputes rank and threshold. The latter deliberately supplies independent protected angles, then assumes their trust. These observations identify the toy state but leave no residual redundancy. Authentication, device security and the selection of trustworthy sensors are outside this demonstration. The same supplied protected evidence goes to fixed and agentic supervision.}
\end{frame}

\begin{frame}[t]{Cyber code 7/12: turn an estimated state into a constrained injection}
\deepdive
\checkpoint{C3}{Read the scalar DC redispatch that uses either estimate.}
\implcode{standard_formulations}{397}{404}
\footnotesize
\textbf{Control meaning:} \texttt{direction} adds active power at node 2, balanced at the source. \texttt{sensitivity} gives the four branch-flow changes per unit injection. The scalar command is bounded between 0 and 0.5 p.u.; squared effort is minimized.
\par\smallskip
\textbf{Constraint meaning:} predicted branch flows must stay below limits $(0.35,0.18,0.35,0.35)$ p.u. with a 0.01 p.u. planning margin. That margin is an engineering choice, not a calibrated uncertainty guarantee. This is conventional optimization, not language-generated control.
\sources{Source: \texttt{standard\_formulations.py:397--404}; control record: \texttt{cyber.trusted\_recovery.control}. State-estimation context: \href{https://doi.org/10.1109/TPAS.1970.292678}{Schweppe and Wildes (1970)}.}
\note{Optional code detail. The four-by-three matrix F converts estimated angles to flows. B is the reduced DC nodal injection matrix; solving with B converts a node-2 injection into angle changes, and multiplication by F gives line-flow sensitivity. SLSQP solves a one-variable convex quadratic problem with linear inequalities and simple bounds in this example. All powers share one fixed toy-model per-unit base. No voltage magnitudes, inverter reactive capability or distribution-feeder unbalance are represented.}
\end{frame}

\begin{frame}[t]{Cyber code 8/12: check predicted and actual toy consequences separately}
\deepdive
\checkpoint{C3}{A returned command is compared with a common synthetic truth.}
\implcode{standard_formulations}{407}{416}
\footnotesize
\textbf{Protected-estimate result:} injection 0.121693667 p.u.; predicted flow 1--2 = 0.170000000 p.u.; actual toy flow = 0.170460762 p.u. All four actual branch limits pass. The difference reflects protected measurement noise.
\par\smallskip
\textbf{Fair evaluation:} the stealth-derived zero injection is evaluated with the same physical truth and leaves one overload. This improvement comes from additional evidence plus conventional redispatch. A model checking its own corrupted estimate would conceal the failure.
\sources{Source: \texttt{standard\_formulations.py:407--416}; actual values: \texttt{results/standard\_formulations.json}. Threat limitation: \href{https://doi.org/10.1145/1653662.1653666}{Liu et al. (2009)}.}
\note{Optional code detail. The post-solve guard checks finite values, command bounds and predicted constraints. The returned actual_violation_count separately evaluates all branches under known synthetic truth. The standard workflow then requires that count to be zero before declaring its modeled verification PASS. The presence of truth is an experimental evaluation device, not a live operator sensor. This remains a DC demonstration with no full-feeder AC security claim.}
\end{frame}

\begin{frame}[t]{Cyber code 9/12: Llama selects a permitted method, not new trust}
\deepdive
\checkpoint{C4}{Follow the recorded Planner response into the application scope gate.}
\implcode{standard_agent_workflow}{113}{120}
\footnotesize
\textbf{Actual call 7:} \texttt{WLS\_TRUSTED\_REDISPATCH}. Returned reason, verbatim: ``This tool is authorized for the given use case and scope, and it is designed for a specific DC redispatch example with trusted measurements.''
\par\smallskip
The specification already declares the authorized method and trust assumption. A wrong method or \texttt{REVIEW} stops the numerical route. In this record the permitted tool runs; Llama does not authenticate sensors or choose a better estimator.
\sources{Source: \texttt{standard\_agent\_workflow.py:113--120}. Actual reason: \texttt{standard\_stage\_20260921T143826Z\_42d19a61/llama\_calls.json}, call 7. Bounded tool use: \href{https://arxiv.org/abs/2210.03629}{ReAct (2023)}.}
\note{Optional code detail. This is the returned explanation text, not hidden chain-of-thought. The enum chooses a permitted next step; numerical parameters remain fixed. The historical standard-method wrapper allows one numerical attempt and does not explore arbitrary methods or recover from unavailable tools. The fixed supervisor maps the identical structured specification directly to the same tool. Therefore a correct selection here is visible functionality, not evidence of superiority over that comparator.}
\end{frame}

\begin{frame}[t]{Cyber code 10/12: the enum and the explanatory prose disagree}
\deepdive
\checkpoint{T6}{The actual Explanation response must retain its stopping effect.}
\implcode{standard_agent_workflow}{134}{141}
\footnotesize
\textbf{Actual call 8: choice = \texttt{REVIEW}. Returned reason, verbatim:}
\par\smallskip
{\itshape
The provided JSON output indicates that the numerical workflow has completed successfully, and the results are within the expected scope. The 'passed\_modeled\_checks' field is set to True, indicating that the modeled checks have been passed. Therefore, the choice is to REVIEW the results to ensure they are accurate and within the expected range.
\par}
\textbf{Implementation consequence:} the structured choice sets acceptance to false. The passing tool result stays provisional; the final workflow status is \texttt{REVIEW}. Favorable prose cannot override the declared stopping choice.
\sources{Source: \texttt{standard\_agent\_workflow.py:134--141}. Exact response: \texttt{results/standard\_stage\_20260921T143826Z\_42d19a61/llama\_calls.json}, call 8; final result: \texttt{standard\_agent\_workflow.json}.}
\note{Optional code detail. Read the enum before interpreting the prose. The local model acknowledges passed checks but chooses REVIEW for another review; the wrapper faithfully retains that escalation. There was no service or schema error on this call. We do not reclassify the output as a recommendation, rerun it until favorable, or imply that escalation improved numerical safety. This is an observed limitation of the historical bounded explanation interface, distinct from the later continuous-voltage contract revision.}
\end{frame}

\begin{frame}[t]{Cyber code 11/12: report measured latency with its actual boundary}
\deepdive
\checkpoint{C4}{These historical timers and token counts are read from saved calls.}
\implcode{standard_agent_workflow}{60}{67}
\footnotesize
\begin{tabular}{@{}lrrrr@{}}
\toprule Llama call & Client (s) & Server (s) & Prompt tokens & Output tokens \\
\midrule Planning 7 & 2.699153 & 0.646972 & 193 & 46 \\
Explanation 8 & 3.230637 & 1.158535 & 419 & 77 \\
\bottomrule
\end{tabular}
\par\smallskip
\textbf{Whole recorded workflow:} fixed 0.002146 s; agentic 5.934650 s. Client time includes the HTTP request and response/schema handling; Ollama's \texttt{total\_duration} is a narrower server measure. This run did not isolate each numerical stage or measure energy or monetary cost.
\sources{Timer source: \texttt{standard\_agent\_workflow.py:43,60--67}; exact client/server fields: \texttt{standard\_stage\_20260921T143826Z\_42d19a61/llama\_calls.json}; whole-run \texttt{seconds}: \texttt{standard\_agent\_workflow.json}.}
\note{Optional code detail. The timer starts immediately before the try block at line43 and stops at line62 after HTTP response decoding and schema handling; journal writing occurs after this measurement. Both actual cyber calls succeeded, so the displayed failure-path lines were not exercised. Client seconds are Python wall time; server total_duration is reported in nanoseconds and is divided by one billion for this table. prompt_eval_count and eval_count provide the two token columns. Their sums are612 prompt and123 generated tokens, not a dollar bill. Whole workflow timing includes its numerical work and agent calls but is a single local observation, not a latency distribution. Do not label the fixed workflow's2.146ms as isolated WLS time or assign unmeasured costs to individual stages.}
\end{frame}

\begin{frame}[t]{Cyber code 12/12: preserve the baseline and hand the record to a person}
\deepdive
\checkpoint{C4}{The stored identity binds the specification, result and stage responses.}
\implcode{standard_agent_workflow}{149}{156}
\footnotesize
\begin{tabular}{@{}lrrl@{}}
\toprule Same protected-data tool & Injection (p.u.) & Actual flow 1--2 (p.u.) & Status \\
\midrule Fixed supervision & 0.12169 & 0.17046 & \texttt{MODEL\_RECOMMEND} \\
Llama supervision & 0.12169 & 0.17046 & \texttt{REVIEW} \\
\bottomrule
\end{tabular}
\par\smallskip
\textbf{Operator questions:} Where did protected evidence come from? Which limits and model were checked? Why did the model request review? Numerical recovery is identical; Llama changes the handoff. A stored hash is an integrity reference, not authentication or authority to operate equipment.
\sources{Source: \texttt{standard\_agent\_workflow.py:149--156}; exact paired records: \texttt{results/standard\_agent\_workflow.json}, cyber fixed/agentic. Human authority: \href{https://doi.org/10.1109/3468.844354}{Parasuraman et al. (2000)}.}
\note{Optional code detail. The first excerpt line completes numerical-failure handling; lines150 onward show the universal pending-human-acknowledgment record and final disposition. A numerical candidate may be retained in summary for explanation even when final status is REVIEW. Neither route operates equipment. The proper comparison preserves identical protected measurements, DC model, optimization and checker. Trusted-data conventional recovery fixes the toy overload, while this Llama record adds latency and an escalation without improving the computed action. An operator may request clarified intent, justified trust or a wider model, rather than accepting favorable numerical text alone.}
\end{frame}
\endgroup


\begin{frame}[t]{Cyber response: what improved, and compared with what?}
\deepdive
\small
\begin{tabular}{@{}p{.37\linewidth}rrl@{}}
\toprule Stealth-case route & Injection $u$ & Actual flow 1--2 & Handoff \\
\midrule
Residual test on corrupted data & 0 & 0.24000 & No alarm \\
Trusted WLS/control + fixed rules & 0.12169 & 0.17046 & Recommend \\
Same trusted tool + Llama & 0.12169 & 0.17046 & Review \\
\bottomrule
\end{tabular}
\par\smallskip\scriptsize Values are per unit; branch limit 0.18. Llama's row is a provisional tool result, evaluated using synthetic truth. Its final status is REVIEW; no action is issued.
\normalsize
\takeaway{Trusted measurements and conventional redispatch remove this toy overload. Adding Llama changes the handoff to REVIEW; the numerical detector and control results stay the same.}
\small\textbf{Operator handoff:} confirm data trust and model scope before acknowledging the review record.
\sources{Numerical baseline and control: \texttt{results/standard\_formulations.json}; actual agent response: \texttt{standard\_agent\_workflow.json}. Stealth attacks: \href{https://doi.org/10.1145/1653662.1653666}{Liu et al. (2009)}.}
\note{162--165 minutes. Close the comparison in two steps. Conventional residual-only estimation fails against this model-consistent attack. Providing trusted observable measurements and a conventional redispatch tool changes the action and removes the overload under known synthetic truth. The same engineering tool yields the same result under fixed and agentic supervision, while the recorded language explanation requests REVIEW. This is an observed escalation rather than evidence of superior detection or safer physical control. The wrapper receives the trusted branch explicitly; Llama neither authenticates nor discovers the trusted sensors. Carry the action, checks, trust condition and model scope into the operator review.}
\end{frame}
